\usepackage[utf8]{inputenc}
\usepackage[italian]{babel}
\usepackage[top=2.5cm, bottom=2.5cm, left=2.5cm, right=2.5cm]{geometry}
\usepackage{graphicx}
\usepackage{tabularx}
\usepackage{booktabs} 
\usepackage{xcolor}
\usepackage{fancyhdr}
\usepackage{hyperref}
\usepackage{eurosym}
\usepackage{textcomp}
\usepackage{xltabular}
\usepackage{longtable}
\usepackage{colortbl}
\usepackage{enumitem}
\usepackage{float}
\usepackage{array}
\usepackage{multirow}
\usepackage{amsmath}
\usepackage{comment}
\usepackage{diagbox}
\usepackage{verbatim}

\usepackage{pgfgantt}     % Gantt
\usepackage{adjustbox}    % scala il Gantt alla larghezza pagina
\usepackage{pdflscape}    % opzionale: ruota la pagina in landscape per il Gantt

\newcolumntype{M}[1]{>{\centering\arraybackslash}m{#1}}
\newcolumntype{Y}{>{\raggedright\arraybackslash}X}

% === PROFONDITA INDICE ===
\setcounter{secnumdepth}{4}
\setcounter{tocdepth}{4}

% === CONFIGURAZIONE COLORI E LINK ===
\definecolor{primarycolor}{RGB}{20, 40, 75}

\hypersetup{
	colorlinks=true,
	linkcolor=primarycolor,
	urlcolor=primarycolor
}

% === LINK INTERNI SOTTOLINEATI E COLORATI ===
\usepackage{soul}

\makeatletter

% Patch per \hyperref
\let\oldhyperref\hyperref
\renewcommand{\hyperref}[2]{%
  \textcolor{primarycolor}{\oldhyperref[#1]{\ul{#2}}}%
}

% Patch per \hyperlink 
\let\oldhyperlink\hyperlink
\renewcommand{\hyperlink}[2]{%
  \textcolor{primarycolor}{\oldhyperlink{#1}{\ul{#2}}}%
}

\makeatother

% === IMPOSTAZIONI PARAGRAFO ===
\setlength{\parindent}{0pt}
\setlength{\parskip}{0.5em}


% === COMANDI PERSONALIZZATI ===
% Comando per link generici (blu e sottolineati)
\newcommand{\link}[2]{\href{#1}{\underline{#2}}}

% Comando per link automatici alle issue GitHub
\newcommand{\ghissue}[1]{\href{https://github.com/Synergo-Unit-UniPD/Second-Brain/issues/#1}{\underline{\##1}}}

%comando per link spezzati su più righe
\usepackage{xurl}
\usepackage{ragged2e}
\usepackage{soul}

\newcommand{\doclink}[2]{%
  \href{#1}{%
    \RaggedRight
    \sloppy
    \ul{#2}%
  }%
}